\documentclass[10pt,a4paper]{article}
\usepackage[margin=24mm,headheight=14pt]{geometry}
\usepackage{fontspec}
\setmainfont{texgyrepagella-regular.otf}[BoldFont=texgyrepagella-bold.otf,ItalicFont=texgyrepagella-italic.otf,BoldItalicFont=texgyrepagella-bolditalic.otf]
\setsansfont{texgyreheros-regular.otf}[BoldFont=texgyreheros-bold.otf]
\setmonofont{texgyrecursor-regular.otf}[Scale=0.85]
\usepackage{microtype}
\usepackage{xcolor}
\usepackage{xurl}
\usepackage{hyperref}
\usepackage{bookmark}
\usepackage{enumitem}
\usepackage{fvextra}
\usepackage{titlesec}
\usepackage{fancyhdr}
\definecolor{navy}{HTML}{172A46}
\definecolor{teal}{HTML}{126E78}
\hypersetup{colorlinks=true,linkcolor=navy,urlcolor=teal,pdftitle={UMI Training Data and Model Licensing Policy},pdfauthor={UMI Inc.},pdfsubject={Conditional model licensing from C7 onward; policy revision October 28, 2026}}
\setlength{\parindent}{0pt}
\setlength{\parskip}{5pt plus 1pt}
\setlength{\emergencystretch}{3em}
\setcounter{secnumdepth}{0}
\setlist{leftmargin=1.6em,itemsep=3pt,topsep=3pt}
\titleformat{\section}{\Large\sffamily\bfseries\color{navy}}{}{0pt}{}
\titleformat{\subsection}{\normalsize\sffamily\bfseries\color{navy}}{}{0pt}{}
\titlespacing*{\section}{0pt}{15pt}{7pt}
\titlespacing*{\subsection}{0pt}{10pt}{3pt}
\providecommand{\tightlist}{\setlength{\itemsep}{0pt}\setlength{\parskip}{0pt}}
\DefineVerbatimEnvironment{verbatim}{Verbatim}{fontsize=\footnotesize,breaklines=true,breakanywhere=true,breaksymbolleft={},frame=single,framesep=3mm}
\pagestyle{fancy}
\fancyhf{}
\fancyhead[L]{\small\sffamily UMI / TRAINING DATA AND MODEL LICENSING}
\fancyhead[R]{\small\sffamily 28 OCT 2026}
\fancyfoot[C]{\small\thepage}
\renewcommand{\headrulewidth}{0.3pt}
\begin{document}
{\huge\sffamily\bfseries\color{navy} UMI Training Data and\\[3pt] Model Licensing Policy\par}
\vspace{4pt}
Policy updated October 28, 2026 · UMI Inc.

Legal assessment and source records dated October 22, 2026.

From C7 onward, UMI uses two model licensing tracks. A model trained
with permitted NonCommercial data, including How2Sign, or derived from
an NC checkpoint must declare that provenance and use CC BY-NC-SA 4.0
with fresh model-bound NC consent. A model whose training data and
inherited components support its existing permissive release keeps the
original contribution terms and declared license. A mixed model that
uses a permitted NC training source follows the NC track.

\textbf{Already submitted to C7? No changes are required because of this
update.} Existing submissions under the original terms do not need
relicensing, re-signing or resubmission. A miner may optionally submit a
new or replacement NC-trained model while intake is open, using the
correct declared license and fresh consent. This update does not
retroactively rewrite accepted artifacts or excuse a false original
rights declaration.

The same two-track approach applies to C8 and later cohorts. NC training or
an inherited NC checkpoint selects the NC track; declaring an NC source
used only for validation or testing does not by itself select that track
or authorize that evaluation use. Retain the applicable source
permissions in either case. Both tracks use the same evaluation,
resource limits, baseline/grace rule and reward formula; a license label
provides no extra score or allocation. Existing content and
score-duplicate rules apply across both tracks. The existing baseline
and earlier certified models retain their own licenses. This revision
incorporates UMI\textquotesingle s conditional policy decision into the
October 22 review; the legal authorities and source records below retain
their stated assessment date.

\section{1. The policy in brief}\label{1-the-policy-in-brief}

The applicable signed contribution terms and authenticated model-license
adoption select the track for each model. The rules below describe the
conditional policy; publishing this review does not change a signed
submission, terms file, cohort clock or historical result.

\begin{enumerate}
\def\labelenumi{\arabic{enumi}.}
\tightlist
\item
  \textbf{Both tracks publish their models.} Certified weights are
  published by UMI on the UMI model repository for anyone to download.
  There is no private or exclusive track.
\item
  \textbf{Choose the license from the model\textquotesingle s actual
  sources and lineage.} On the NC track, weights, the portable bundle
  and adapter files carry CC BY-NC-SA 4.0. On the permissive track,
  retain the existing compatible license and original contribution
  terms. UMI-authored runtime code stays Apache-2.0; your own inference
  code may be Apache-2.0 or MIT. A cohort containing an NC model does
  not make other entries NC.
\item
  \textbf{You may train on permitted NonCommercial datasets through the
  NC track.} Datasets under CC BY-NC 4.0 or CC BY-NC-SA 4.0 with no
  additional contractual restrictions are permitted, including How2Sign.
  Section 3 lists the license classes and the rulings for named
  datasets.
\item
  \textbf{You may not train on datasets with bespoke research or access
  contracts} (for example the Microsoft Research License Terms that
  govern ASL STEM Wiki) unless UMI has the rights
  holder\textquotesingle s written permission on file. Section 9 lists
  the permissions UMI is pursuing.
\item
  \textbf{Every submission carries a provenance declaration and a rights
  warranty} listing each dataset by name, version, URL, license and
  download route (section 5).
\item
  \textbf{Every bundle attributes its sources} in a NOTICE file in the
  form section 4 specifies, and reproduces each
  dataset\textquotesingle s responsible-use requests.
\item
  \textbf{Awards are retroactive.} Emissions for a cohort are set only
  after its results are certified and run for a bounded window until the
  next cohort is certified. Public release under the selected model
  license is a condition of the award; on the NC track that means CC
  BY-NC-SA 4.0.
\item
  \textbf{NC rights do not authorize commercial deployment.} UMI does
  not run inference on NC-track models for money, bundle, link, default
  to or document them in any UMI app, or use them or their outputs to
  build a product model. Permissive-track models remain governed by
  their own licenses and original grants; those rights are not narrowed
  by another entry\textquotesingle s NC license. A
  user\textquotesingle s use remains subject to the selected
  model\textquotesingle s license.
\item
  \textbf{You delete what you must delete.} Where a
  dataset\textquotesingle s terms require destruction after use, you
  destroy your copy and all derived files after training and attest to
  it.
\item
  \textbf{No faces in publicity.} Dataset videos never appear in demos
  or marketing without blurring where the dataset asks for it, and never
  where its terms forbid it.
\end{enumerate}

\section{2. Why NC-track models are non-commercial research
models}\label{2-why-nc-track-models-are-non-commercial-research-models}

For the NC track, UMI\textquotesingle s cohorts convert paid effort into
a public good in the same way a university converts salaried effort into
publications, and the NonCommercial test is applied to that use, not to
the people who perform it or to the money that funds it. The rules in
section 1 are the facts that make this true, and each of them maps to a
reason.

\textbf{The test looks at the use.} The Creative Commons NonCommercial
element permits any use that is "not primarily intended for or directed
towards commercial advantage or monetary compensation." Creative Commons
applies this to the use, not the user: a nonprofit can breach it and a
company can comply with it, and the courts that have construed it
(\emph{Great Minds v. FedEx Office}, 2d Cir. 2018; \emph{Great Minds v.
Office Depot}, 9th Cir. 2019) measure the licensee\textquotesingle s own
use and ignore paid intermediaries. The salary of the person doing the
work, the revenue of the institution and the motives of the participants
are all outside the test. University researchers are paid; universities
sell tickets and charge tuition; their research use of NC data is
treated as non-commercial everywhere. The same test cannot invoke pay,
revenue or motive against UMI.

\textbf{The use is the production of a public research model.} The
required output on the NC track is a model released to the public under
CC BY-NC-SA 4.0, certified against a benchmark, with results published.
That is research output in the ordinary machine-learning sense. Public
release is the outcome that research licensors want most, and it is more
than much university research delivers, since much of that is never
disclosed.

\textbf{The award is recognition, not purchase.} Emissions for a cohort
are set only after certification, for a bounded window, for a model that
already exists in public. No buyer acquires the model, nothing is
exclusive and nothing is transferred. This is retroactive public-goods
funding, an established grant form. Grants are awarded on expected
quality, prizes and renewals on delivered quality, so a payment that
tracks quality is not what distinguishes a grant from a sale; what
distinguishes them is whether the funder acquires and exploits the
output or recognises a contribution to the public. UMI acquires nothing.

\textbf{The NC-track output stays non-commercial.} CC BY-NC-SA 4.0 on
the weights is the downstream condition the NC license itself imposes on
shared adaptations, and Creative Commons\textquotesingle{} own guidance
on AI training treats a model shared non-commercially as compliant. UMI
does not deploy NC-track models commercially, does not put them in its
app and does not build product models from them, so the public good
remains a public good after certification.

\textbf{The licensor\textquotesingle s intent is honoured.} Research
licensors use NC and research-only language to prevent appropriation,
secrecy and commercial products. A public NC-SA model involves none of
those. How2Sign\textquotesingle s signers consented to public
availability for research, and a public research model is what they get;
the How2Sign site\textquotesingle s "research purposes only" sentence
sits on top of a CC BY-NC 4.0 grant and describes exactly this outcome.

\textbf{Deployment is the user\textquotesingle s act.} For NC models,
the UMI app acts as a general-purpose on-device runtime that loads a
model a user independently supplies; UMI does not supply an NC model or
a pointer to one. A user who loads an NC model for personal,
non-commercial use receives the license automatically and makes a
licensed use; where no one in the chain infringes, there is nothing for
secondary liability to attach to, and a general-purpose runtime with
substantial lawful uses is itself lawful. UMI\textquotesingle s model
publication on this track is non-commercial, and the app never points at
an NC-track model. This does not revoke the original rights in
permissive models.

These reasons are only as good as the facts, which is why section 1 is
written as rules and why section 6 keeps UMI\textquotesingle s own
conduct inside them.

\section{3. Which datasets you may train
on}\label{3-which-datasets-you-may-train-on}

A dataset is eligible when the terms attached to the acquired copy
permit the intended training use and model release, or an applicable
recorded permission covers them. Permitted NC training sources select
the NC track and its public CC BY-NC-SA 4.0 release. Sources and
inherited components that support the existing permissive release retain
that route. Additional access contracts must be disclosed and cleared;
an NC label does not override them. Classify by the terms actually
attached to the copy you download, not by what a paper or a mirror says.

\subsection{CC0, public domain}\label{cc0-public-domain}

\textbf{Eligible?:} Yes

\textbf{Why:} No conditions

\subsection{CC BY 4.0, CC BY 3.0}\label{cc-by-40-cc-by-30}

\textbf{Eligible?:} Yes

\textbf{Why:} Attribution only; carried in NOTICE

\subsection{CC BY-SA 4.0, CC BY-SA 3.0 (FLEURS-ASL, FLORES, Wikipedia
text)}\label{cc-by-sa-40-cc-by-sa-30-fleurs-asl-flores-wikipedia-text}

\textbf{Eligible?:} Yes

\textbf{Why:} Attribution and ShareAlike; UMI\textquotesingle s recorded
position is that weights are not an adaptation of the dataset. Preserve
applicable component licenses. Permissive-only models retain their
compatible release; mixed models containing permitted NC training
sources use NC-SA. This is not a general claim that BY-SA adaptations
may be relicensed NC-SA; section 9 records the requested Google
confirmation

\subsection{CC BY-NC 4.0, CC BY-NC-SA 4.0
(How2Sign)}\label{cc-by-nc-40-cc-by-nc-sa-40-how2sign}

\textbf{Eligible?:} Yes

\textbf{Why:} Non-commercial research use and NC-SA release match the
license\textquotesingle s own conditions

\subsection{Apache-2.0, MIT, BSD, ODC-By,
CDLA-Permissive}\label{apache-20-mit-bsd-odc-by-cdla-permissive}

\textbf{Eligible?:} Yes

\textbf{Why:} Attribution and notice only

\subsection{CC BY-ND, CC BY-NC-ND, any no-derivatives
term}\label{cc-by-nd-cc-by-nc-nd-any-no-derivatives-term}

\textbf{Eligible?:} No

\textbf{Why:} Training produces an adaptation on the
licensor\textquotesingle s own reading

\subsection{Bespoke research or access agreements (Microsoft Research
License Terms, ELRA end-user licences, click-through academic
agreements)}\label{bespoke-research-or-access-agreements-microsoft-research-license-terms-elra-end-user-licences-click-through-academic-agreements}

\textbf{Eligible?:} No, unless UMI holds written permission

\textbf{Why:} These are contracts with terms beyond NC, typically no
distribution of modifications, destruction after use, named-institution
or named-user limits

\subsection{Scraped video, YouTube, social platforms, unlicensed
mirrors}\label{scraped-video-youtube-social-platforms-unlicensed-mirrors}

\textbf{Eligible?:} No

\textbf{Why:} No license; platform terms prohibit it; signer consent
absent

\subsection{UMI consented data (collected through the UMI app under the
UMI contributor
consent)}\label{umi-consented-data-collected-through-the-umi-app-under-the-umi-contributor-consent}

\textbf{Eligible?:} Yes

\textbf{Why:} Apply the actual contributor consent and license to the
selected model track; do not infer broader rights

\textbf{Rulings for named datasets.} These are the datasets miners have
asked about. Rights-holder permissions are recorded in section 9. The
specific ChicagoFSWild rulings below record the separate operator
approval; they do not invent publisher license terms.

\subsection{\texorpdfstring{\href{https://www.kaggle.com/datasets/googleai/fleurs-asl}{FLEURS-ASL}
(Google Research, CC BY-SA
4.0)}{FLEURS-ASL (Google Research, CC BY-SA 4.0)}}\label{fleurs-asl-google-research-cc-by-sa-40}

\textbf{Ruling:} Allowed

\textbf{Conditions:} Attribute FLEURS-ASL and FLORES in NOTICE;
reproduce the Responsible Use requests (blur faces in publicity, no
reidentification, no deepfakes)

\subsection{\texorpdfstring{\href{https://how2sign.github.io/}{How2Sign}
(UPC, BSC, CMU, Gallaudet; CC BY-NC
4.0)}{How2Sign (UPC, BSC, CMU, Gallaudet; CC BY-NC 4.0)}}\label{how2sign-upc-bsc-cmu-gallaudet-cc-by-nc-40}

\textbf{Ruling:} Allowed with conditions

\textbf{Conditions:} Attribute the How2Sign consortium and cite the CVPR
2021 paper; release under CC BY-NC-SA 4.0; do not redistribute the
videos, clips, landmarks or English sentences; do not train the How2
source transcripts into a model intended for any other track

\subsection{\texorpdfstring{\href{https://www.microsoft.com/en-us/research/project/asl-stem-wiki/}{ASL
STEM Wiki} (Microsoft, Microsoft Research License
Terms)}{ASL STEM Wiki (Microsoft, Microsoft Research License Terms)}}\label{asl-stem-wiki-microsoft-microsoft-research-license-terms}

\textbf{Ruling:} Pending permission

\textbf{Conditions:} Not permitted until Microsoft\textquotesingle s
written confirmation is recorded in section 9; the terms forbid
distributing modifications and require destruction after use, which the
NC track does not satisfy on its own

\subsection{2M-Flores-ASL, FSboard, Taskmaster-1 (UMI reference-model
lineage)}\label{2m-flores-asl-fsboard-taskmaster-1-umi-reference-model-lineage}

\textbf{Ruling:} Allowed

\textbf{Conditions:} Attribute as the reference model\textquotesingle s
NOTICE does; each is CC BY or CC BY-SA

\subsection{UMI consented data}\label{umi-consented-data}

\textbf{Ruling:} Allowed

\textbf{Conditions:} Attribute as "UMI contributors"; no per-contributor
identification; retain the actual consent conditions

\subsection{ChicagoFSWild and
ChicagoFSWild+}\label{chicagofswild-and-chicagofswild}

\textbf{Ruling:} Operator-approved train splits only

\textbf{Conditions:} Disclose use in PROVENANCE; never train on
dev/test. Retain the actual publisher terms, including absence of an
explicit license where applicable. This specific approval does not
authorize other unlicensed sources or override a separate access
contract

A dataset not listed here is eligible if it falls in a "Yes" class above
and you record its terms in your provenance declaration. If you are
unsure which class applies, ask before training; rulings are added to
this table as they are made.

\section{4. How to license and label your
bundle}\label{4-how-to-license-and-label-your-bundle}

The bundle carries one license per component and a NOTICE file that lets
anyone trace every source. A bundle label never widens the rights in the
training data; it records the rights that exist.

\subsection{NC-track weights, adapters, tokenizer and label
vocabularies, portable
bundle}\label{nc-track-weights-adapters-tokenizer-and-label-vocabularies-portable-bundle}

\textbf{License:} CC BY-NC-SA 4.0

\textbf{What the file must say:} \texttt{LICENSE-WEIGHTS} containing the
CC BY-NC-SA 4.0 legal code or its URL; declared model identifier
\texttt{CC-BY-NC-SA-4.0}

\subsection{Permissive-track weights and portable
bundle}\label{permissive-track-weights-and-portable-bundle}

\textbf{License:} Existing compatible declared license

\textbf{What the file must say:} Retain the original component licenses
and notices; do not replace them with NC-SA merely because the cohort
also accepts NC models

\subsection{Your inference code and
configs}\label{your-inference-code-and-configs}

\textbf{License:} Apache-2.0 or MIT

\textbf{What the file must say:} \texttt{LICENSE} for the code

\subsection{UMI runtime and extractor code you
redistribute}\label{umi-runtime-and-extractor-code-you-redistribute}

\textbf{License:} Apache-2.0 (unchanged)

\textbf{What the file must say:} UMI\textquotesingle s own
\texttt{LICENSE} and \texttt{THIRD\_PARTY\_NOTICES} retained

\subsection{Model card}\label{model-card}

\textbf{License:} CC BY 4.0

\textbf{What the file must say:} \texttt{MODEL\_CARD.md}

\subsection{NOTICE}\label{notice}

\textbf{License:} n/a

\textbf{What the file must say:} \texttt{NOTICE} as specified below

\textbf{The NOTICE file.} One entry per dataset, in this form, and one
entry for any base checkpoint you fine-tuned:

\begin{verbatim}
SOURCE: How2Sign v1.0
HOLDER: How2Sign consortium (Universitat Politècnica de Catalunya, Barcelona Supercomputing Center, Carnegie Mellon University, Gallaudet University)
LICENSE: CC BY-NC 4.0, https://creativecommons.org/licenses/by-nc/4.0/
URL: https://how2sign.github.io/
CITE: Duarte et al., "How2Sign: A Large-scale Multimodal Dataset for Continuous American Sign Language," CVPR 2021
USE: Training data for this model. The dataset itself is not redistributed. The model was trained in part on this source and is released under CC BY-NC-SA 4.0 under UMI's NC-track policy.
REQUESTS: Research use; no redistribution of videos; no identification of signers.
\end{verbatim}

For FLEURS-ASL the entry names Google Research, CC BY-SA 4.0, the Kaggle
URL, the NAACL 2025 paper, and reproduces the Responsible Use paragraph
(blur faces when publicising examples; no reidentification; no
deepfakes). For FLORES text used through FLEURS-ASL, name FLORES and CC
BY-SA 4.0. For Wikipedia-derived text, name Wikipedia and CC BY-SA. For
UMI consented data, name "UMI contributors" and the UMI contributor
consent.

\textbf{The model card} states the task, the architecture in one
paragraph, every contributor-used training and evaluation dataset by
name and version, each component\textquotesingle s license, any
memorisation test result and its method, and the cohort and
certification date when available. NC models additionally state: "This
is a research model released under CC BY-NC-SA 4.0. It is not a product,
not an interpreter, and not for commercial use." Permissive models state
their actual license and intended-use limits instead. This declaration
concerns the contributor\textquotesingle s own sources; it does not
request or disclose UMI\textquotesingle s private challenge provenance.

\textbf{Fine-tuning a certified checkpoint.} Any adaptation of a CC
BY-NC-SA 4.0 cohort checkpoint must be released under CC BY-NC-SA 4.0,
attribute the checkpoint and its cohort, and inherit its NOTICE entries
in full.

\section{5. What to submit with a
model}\label{5-what-to-submit-with-a-model}

For a new or replacement NC submission, include the records below in the
authenticated bundle and accept the exact selected NC terms with fresh
model-bound consent. Use the current structured provenance format and
supported client linked after this list. Existing permissive submissions
retain their original required format and reader; this revision does not
impose a retrospective resubmission gate.

\begin{enumerate}
\def\labelenumi{\arabic{enumi}.}
\tightlist
\item
  \textbf{Provenance declaration.} A \texttt{PROVENANCE.json} listing
  every training, validation and test source and every inherited
  checkpoint: name, version, canonical URL,
  license (SPDX identifier where one exists), the route you downloaded
  it by (direct link, Kaggle, Google Drive, mirror), the date, the
  archive file names with SHA-256 hashes, and any terms you accepted to
  obtain it (account terms of service, click-through, request form). A
  dataset you cannot describe this way is not eligible.
\item
  \textbf{Rights warranty.} A signed statement that each listed dataset
  falls within an applicable permitted license class or named source
  ruling in section 3, or has the required rights-holder permission
  recorded in section 9, that any accepted access terms are accurately
  disclosed and any additional dataset contract has the required
  recorded permission, that no training video, frame, landmark cache,
  caption or sentence list is included in the bundle, and that you grant
  UMI permission to evaluate, preserve, publish and use the exact
  submission as a research baseline under CC BY-NC-SA 4.0 as specified
  in the selected NC supplement. This replaces the original
  commercial-use grant for that model only; permissive models retain
  their original grants.
\item
  \textbf{Memorisation-test reporting.} The original review calls for a
  regurgitation probe using held-out video from each training source,
  with per-source verbatim-match rates and longest matched spans. The
  review does not specify a numeric threshold. UMI must publish the
  actual probe rule before claiming it is an enforced certification
  gate. Report any completed probe and its method honestly; a
  declaration or ordinary benchmark score does not establish that an
  independent probe passed. This revision neither invents a threshold
  nor retroactively changes signed scoring rules.
\item
  \textbf{Bundle inventory.} A file list with hashes, and a statement
  that the bundle contains weights, code, configs, model card and
  notices only.
\item
  \textbf{Destruction attestation}, where a dataset\textquotesingle s
  terms require destruction after use: the date you deleted the dataset
  and all derived files.
\item
  \textbf{NOTICE and model card} as section 4 specifies.
\end{enumerate}

UMI keeps the declaration, warranty, any test output and inventory with
the certified checkpoint\textquotesingle s hash for as long as the
checkpoint is published.

\textbf{Supported submission path.} Use the
\href{https://github.com/Umi-BitSign/umi/blob/main/docs/contributors/models.md}{model
contribution guide},
\href{https://github.com/Umi-BitSign/umi/blob/main/docs/reference/model-training-provenance.md}{structured
provenance format},
\href{https://github.com/Umi-BitSign/umi/blob/main/docs/MODEL_CONTRIBUTION_TERMS_NC_V1.md}{NC
supplement}, and
\href{https://github.com/Umi-BitSign/umi/blob/main/docs/reference/commands.md\#optional-non-commercial-model-terms}{NC
client commands}. Update the supported model client before making an NC
submission; it does not update itself in the background. Authenticate
the cohort\textquotesingle s selected model-license route, prepare the
bundle, accept the exact terms hash, and sign both the submission and
cohort consent. Retain all original artifacts and receipts. No action is
required for an unchanged existing C7 submission.

\section{6. Certification, awards, and what UMI does with certified
models}\label{6-certification-awards-and-what-umi-does-with-certified-models}

Both tracks recognize a public model under the same cohort scoring and
allocation rules. The NC track additionally restricts
UMI\textquotesingle s use to the research-model rights below; the
permissive track retains its original contribution grant and compatible
license.

\textbf{Certification.} A cohort closes on its published date.
Validators use the selected signed scoring and artifact-review rules for
each complete submission, with sealed benchmark references and the
applicable source declaration. Certified models are published with their
actual license and retained notices. NC models carry the non-commercial
condition; permissive models retain their original license. The
memorisation-probe enforcement limit in section 5 remains explicit; this
document does not claim a deployed probe that has no published rule.

\textbf{Awards.} Incentive weights for a cohort go live only after its
results are certified and remain live until the next cohort is
certified. The award is for the certified public model; it is not a fee
for serving inference, and NC models are not served by UMI or validators
as a paid service. Public release under the selected model license and
the applicable signed contribution terms is a condition of the award. On
the NC track, this requires CC BY-NC-SA 4.0. A later cohort does not
relicense a previously certified model; any withdrawal or violation is
handled under the applicable signed terms.

\textbf{What UMI does with certified models.} UMI stores, evaluates and
publishes models, uses them as baselines and reference points, and cites
research results under each model\textquotesingle s applicable grant.
For NC models, UMI does not run paid inference, include or link them in
an app, set them as defaults, document them for an app, or use them or
their outputs to build or bootstrap a product model. The general-purpose
runtime may load an NC model a user independently supplies, subject to
that model\textquotesingle s license; UMI supplies no NC model or
pointer through the app. Permissive models retain the uses supported by
their existing licenses and original grants. Referencing a model as a
baseline never changes its license.

\textbf{Publicity.} Dataset videos and frames appear in UMI materials
only where the dataset\textquotesingle s terms and the
signers\textquotesingle{} consent allow it, with faces blurred where the
dataset asks. NC-track models are described publicly as research models,
never as product or app models. Describe permissive models consistently
with their actual license and permitted use.

\textbf{Takedown.} A rights holder or a signer who objects to a
certified model may write to UMI; UMI will unpublish the checkpoint and
stop using it as a baseline while the objection is reviewed, and will
record the outcome in section 9.

\section{7. Dataset terms on record}\label{7-dataset-terms-on-record}

The table records the license text, access terms, signer consent and
provenance of each dataset named in section 3, as read from the live
pages and metadata. It is the reference for the rulings above and is
updated when a rights holder\textquotesingle s reply is recorded in
section 9.

\subsection{\texorpdfstring{\href{https://www.microsoft.com/en-us/research/project/asl-stem-wiki/}{ASL
STEM Wiki}}{ASL STEM Wiki}}\label{asl-stem-wiki}

\textbf{Rights holder and version:} Microsoft Corporation; v1.0,
published 15 Oct 2024, ASL\_STEM\_Wiki.zip, 187.4 GB,
\href{https://www.microsoft.com/en-us/download/details.aspx?id=106253}{Download
Center}

\textbf{License text in force:}
\href{https://www.microsoft.com/en-us/research/project/asl-stem-wiki/dataset-license/}{Microsoft
Research License Terms}: use "solely for non-commercial, non-revenue
generating, research purposes" (s.1); data may be used and modified only
"consistent with the consent under which the data was provided" and "you
may not distribute the data or your modifications to the data" (s.1(c));
no sharing, publishing, distributing or hosting of the Materials
(s.2(e)); personal data used only for consented purposes, kept in strict
confidence, never shared, and destroyed "immediately upon the completion
of your research" (s.3); results may be published if no material or
substantial portion of the Materials is included (s.5); disputes by
individual arbitration, King County WA venue (ss.9, 11)

\textbf{Access and acquisition terms:} Direct zip download, no
click-through on the download page; the project page points to the
license tab and says commercial use requires contacting
ASL\_\href{mailto:Citizen@microsoft.com}{\nolinkurl{Citizen@microsoft.com}}

\textbf{Signer consent on record:}
\href{https://www.microsoft.com/en-us/research/project/asl-stem-wiki/consent-form/}{Consent
form}: data used "primarily to perform research"; "may eventually be
used to develop and improve our commercial products"
(Microsoft\textquotesingle s); public dataset release "to help advance
research"; license of image and likeness to Microsoft "for the research
project above and in any related services"; collected under IRB
approval; 37 paid professional interpreters, faces visible

\textbf{Provenance notes:} 64,266 videos of 254 English Wikipedia STEM
articles; the English text layer is Wikipedia content under CC BY-SA 3.0
and GFDL
(\href{https://www.microsoft.com/en-us/research/project/asl-stem-wiki/datasheet/}{datasheet});
datasheet says the data "should not be used to develop technology that
purports to replace sign language interpreters"

\subsection{\texorpdfstring{\href{https://how2sign.github.io/}{How2Sign}}{How2Sign}}\label{how2sign}

\textbf{Rights holder and version:} How2Sign consortium (UPC, BSC, CMU,
Gallaudet; "©Copyright How2Sign"); v1.0; ELRA catalogue entry ELRA-S0416
available from 29 Sep 2022; CORA.RDR deposit
\href{https://b2find.eudat.eu/dataset/99ba11b8-9428-5177-9e25-ccccc296e061}{doi:10.34810/data33}
v1.0, 2024

\textbf{License text in force:}
\href{https://creativecommons.org/licenses/by-nc/4.0/}{CC BY-NC 4.0}:
reproduction, sharing and Adapted Material "for NonCommercial purposes
only" (s.2(a)(1)); NonCommercial means "not primarily intended for or
directed towards commercial advantage or monetary compensation"
(s.1(i)); attribution per s.3(a); rights terminate automatically on
breach (s.6(a)); no restriction on uses lawful without permission
(s.8(a)). Site adds: "publicly available for research purposes only" and
"you may NOT use the material for commercial purposes"

\textbf{Access and acquisition terms:} Google Drive links plus a GitHub
download script; no click-through or request form on the site; ELRA
lists the same CC-BY-NC-4.0 licence for academic and commercial users at
EUR 0 (no commercial licence offered); the CORA.RDR deposit carries
"custom dataset terms" behind a bot wall, not yet retrieved (section 9)

\textbf{Signer consent on record:}
\href{https://arxiv.org/pdf/2008.08143v2}{Paper s.3.5}: 11 signers (4
Deaf, 2 hard of hearing, 5 hearing) signed a consent form "agreeing on
being recorded and making their data publicly available for research
purposes"; CMU IRB; faces visible

\textbf{Provenance notes:} \textasciitilde80 hours of ASL translations
of How2 instructional videos; English sentences come from the
uploader-provided subtitle track of the How2 source videos; the site
says "for copyright reasons, we are not allowed to redistribute the How2
data"

\subsection{\texorpdfstring{\href{https://www.kaggle.com/datasets/googleai/fleurs-asl}{FLEURS-ASL}}{FLEURS-ASL}}\label{fleurs-asl}

\textbf{Rights holder and version:} Google Research; Kaggle dataset
published 13 Aug 2024, last modified 29 Aug 2024, archive 24.79 GB;
paper \href{https://aclanthology.org/2025.naacl-long.314/}{NAACL 2025}

\textbf{License text in force:} CC BY-SA 4.0 (Kaggle Croissant metadata:
license name "CC BY-SA 4.0", creativecommons.org/licenses/by-sa/4.0;
paper footnote: released "under CC BY-SA 4.0"). Attribution per s.3(a);
ShareAlike per s.3(b) applies to Adapted Material that is Shared; s.8(a)
as above

\textbf{Access and acquisition terms:} Kaggle account and Kaggle terms
of service required to download; no dataset-specific click-through.
Description adds non-binding "Responsible Use" requests: blur faces when
publicizing examples, do not reidentify interpreters, do not use
likenesses for deepfakes

\textbf{Signer consent on record:} 5 paid Certified Deaf Interpreters;
"all interpreters gave informed consent for the use of this data in
machine learning" and to publication for ML research by the broader
community; faces unblurred in the release

\textbf{Provenance notes:} 7.49 hours, 1,749 sentences; the English
layer is FLORES dev/devtest text (FLORES+ successor is
\href{https://huggingface.co/datasets/openlanguagedata/flores_plus}{CC
BY-SA 4.0}), which is why Google chose ShareAlike; the
\href{https://huggingface.co/datasets/google/fleurs}{FLEURS} speech card
says CC BY 4.0, an inconsistency that does not affect
FLEURS-ASL\textquotesingle s own license

\section{8. Legal basis}\label{8-legal-basis}

For the NC track, UMI\textquotesingle s recorded legal position is that
cohort training on NonCommercial datasets is a licensed non-commercial
research use; that, independently, the training copies are fair use in
the United States and text-and-data-mining in the EU; that trained
weights are not a copy or adaptation of any dataset; and that deployment
by users in a general-purpose runtime is the users\textquotesingle{} own
licensed act. The position is stated under United States law as of 22
October 2026, with EU and UK notes where the result differs, and rests
on the authorities in the table that follows.

\textbf{I. The NonCommercial test is use-based.} CC BY-NC 4.0 s.1(i)
defines NonCommercial as "not primarily intended for or directed towards
commercial advantage or monetary compensation." Creative
Commons\textquotesingle{} interpretation guidance makes the test turn on
the use rather than the user; \emph{Great Minds v. FedEx Office}, 886
F.3d 91 (2d Cir. 2018), and \emph{Great Minds v. Office Depot} (9th Cir.
2019) measure the licensee\textquotesingle s own use and disregard paid
intermediaries; Creative Commons\textquotesingle{} guidance on AI
training (28 May 2025) states that the license applies only where
copyright permission is required, that NC is satisfied when each
permission-requiring step is non-commercial, and that a model shared
non-commercially complies. Section 8(a) of every CC 4.0 license confirms
that the license imposes no condition on a use that is lawful without
permission.

\textbf{II. NC-track cohort training is a non-commercial research use.}
The use on this track is the reproduction of a dataset to train a model
whose required output is a public model released under CC BY-NC-SA 4.0,
certified against a benchmark, with published results. The award is
retroactive, bounded, and recognises a deliverable that already exists
in public without any transfer or exclusivity, which is the structure of
retroactive public-goods funding rather than a purchase. The
miners\textquotesingle{} motives and compensation are outside the test
for the same reason a researcher\textquotesingle s salary is. The output
is shared on the terms the NC license itself prescribes for adaptations,
and the licensor\textquotesingle s purposes in choosing NC, the
prevention of appropriation, secrecy and commercial products, are served
by the release. How2Sign\textquotesingle s "research purposes only"
sentence sits on top of a CC BY-NC 4.0 grant with no assent mechanism
and describes this outcome.

\textbf{III. The training copies are independently lawful without any
license.} Intermediate, non-expressive copying to build a functional
tool is fair use under \emph{Sega v. Accolade} (9th Cir. 1992),
\emph{Sony v. Connectix} (9th Cir. 2000), \emph{Authors Guild v. Google}
(2d Cir. 2015) and \emph{Google v. Oracle} (2021), applied to model
training in \emph{Bartz v. Anthropic} (N.D. Cal. 23 Jun 2025) and
\emph{Kadrey v. Meta} (N.D. Cal. 25 Jun 2025). The works are thin
(green-screen linguistic recordings; factual transcripts); the model
does not substitute for a benchmark corpus; no commercial license is
offered anywhere for the NC datasets, which the Copyright
Office\textquotesingle s Part 3 report treats as weighing toward fair
use; and a public NC release is nonprofit dissemination. \emph{Thomson
Reuters v. Ross} (3d Cir. 29 Sep 2026) concerned a non-generative tool
competing directly in a market where licensing was available and
expressly reserved generative models. In the EU, Art. 4 of Directive
2019/790 covers the training copies for lawfully accessible content not
reserved by the rights holder, and Creative Commons\textquotesingle{}
statement of 17 Dec 2021 confirms that CC licenses are not Art. 4
reservations.

\textbf{IV. Weights are not a copy or adaptation of the datasets.}
\emph{Getty Images v. Stability AI} {[}2025{]} EWHC 2863 (Ch) holds that
a model which stores no copy of the works is not an infringing copy, and
the Copyright Office confines infringement by weights to retained,
memorised protectable expression. A landmark-based model that emits
short English strings retains no video expression; the proposed
memorisation probe in section 5 is intended to assess this, subject to
its explicit enforcement limits. Neither a benchmark score nor a license
declaration proves that a checkpoint retains no protectable expression.
On this recorded analysis, CC BY-NC-SA 4.0 is UMI\textquotesingle s
policy choice for a model trained with permitted NC sources or derived
from an NC checkpoint, rather than a blanket license for all cohort
models. Permissive models retain their compatible licenses;
FLEURS-ASL\textquotesingle s ShareAlike, which requires a BY-SA adapter
license only for adaptations, is addressed by the confirmation UMI is
seeking from Google Research (section 9).

\textbf{V. Deployment by users is the users\textquotesingle{} licensed
act.} For NC models, the UMI app acts as a general-purpose on-device
runtime with substantial lawful uses and loads a model the user
independently supplies; it ships with no NC model or pointer to one. A
user who loads a CC BY-NC-SA 4.0 model for personal non-commercial use
receives the license under s.2(a)(5) and makes a licensed use; absent
direct infringement there is no secondary liability (\emph{Sony v.
Universal}; \emph{MGM v. Grokster}), and general-purpose runtimes are
lawful (\emph{Sony v. Connectix}; \emph{Sony v. Bleem}).
UMI\textquotesingle s publication of the NC model is non-commercial, and
section 6 keeps UMI pointers to NC models out of the app. Uses of
permissive models remain governed by their original licenses and grants.

\textbf{VI. Bespoke contracts are outside this basis.} Datasets governed
by individual license agreements, such as ASL STEM Wiki under the
Microsoft Research License Terms, impose obligations (no distribution of
modifications, destruction after use, use within the data
subjects\textquotesingle{} consent) that are contractual and do not turn
on the analysis above. They are eligible only on the rights
holder\textquotesingle s written permission.

\textbf{Privacy and publicity.} The datasets show
interpreters\textquotesingle{} faces. A text-output model stores no
likeness and identifies no one; landmark extraction for translation is
not biometric identification; in the EU the processing rests on
legitimate interests under EDPB Opinion 28/2024; dataset videos are used
in publicity only as section 6 allows.

\textbf{Authorities.} Settled rules are the license texts and the cases
construing them; the application of fair use, TDM exceptions and the
status of weights to model training is still being decided, and the
table records where each authority stood on 22 October 2026.

\subsection{\texorpdfstring{\emph{Bartz v.
Anthropic}}{Bartz v. Anthropic}}\label{bartz-v-anthropic}

\textbf{Court and date:} N.D. Cal., 23 Jun 2025 (Alsup)

\textbf{What it decided:} Training on lawfully acquired books is
"exceedingly transformative" fair use; a permanent library of pirated
copies is not

\textbf{Status on 22 Oct 2026:} Class settlement (USD 1.5 bn) finally
approved 20 Jul 2026; fair-use ruling never reviewed on appeal; only fee
appeals pending in the Ninth Circuit

\subsection{\texorpdfstring{\emph{Kadrey v.
Meta}}{Kadrey v. Meta}}\label{kadrey-v-meta}

\textbf{Court and date:} N.D. Cal., 25 Jun 2025 (Chhabria)

\textbf{What it decided:} Training Llama held highly transformative;
plaintiffs failed on market harm; market-dilution theory noted as the
strongest unargued objection

\textbf{Status on 22 Oct 2026:} Distribution claim set for hearing Feb
2027; interlocutory appeal denied 8 Jul 2026

\subsection{\texorpdfstring{\emph{Thomson Reuters v.
Ross}}{Thomson Reuters v. Ross}}\label{thomson-reuters-v-ross}

\textbf{Court and date:} 3d Cir., 29 Sep 2026 (precedential), affirming
D. Del. 11 Feb 2025

\textbf{What it decided:} Fair use rejected for a non-generative
legal-research AI trained on headnotes: commercial, minimally
transformative, harm to the licensing market for AI-training data;
footnote expressly leaves generative AI open

\textbf{Status on 22 Oct 2026:} Rehearing petition deadline extended to
3 Nov 2026; cert petition signalled

\subsection{US Copyright Office, Copyright and AI Part
3}\label{us-copyright-office-copyright-and-ai-part-3}

\textbf{Court and date:} Pre-publication report, 9 May 2025

\textbf{What it decided:} Weights can infringe the reproduction and
derivative-work rights where the model retains substantial protectable
expression; copying weights infringes only with substantial similarity;
commerciality turns on substance, not labels; where a licensing market
exists unlicensed use is disfavored

\textbf{Status on 22 Oct 2026:} Still pre-publication; final version
promised without substantive change

\subsection{\texorpdfstring{\emph{Getty Images v. Stability
AI}}{Getty Images v. Stability AI}}\label{getty-images-v-stability-ai}

\textbf{Court and date:} England and Wales High Court, {[}2025{]} EWHC
2863 (Ch), 4 Nov 2025

\textbf{What it decided:} A trained model that does not store or
reproduce the works is not an "infringing copy" under CDPA s.27;
secondary-infringement claim dismissed; primary training claims
withdrawn because training happened outside the UK

\textbf{Status on 22 Oct 2026:} Permission to appeal the s.27 point
granted Dec 2025; Court of Appeal date not yet listed

\subsection{\texorpdfstring{\emph{GEMA v.
OpenAI}}{GEMA v. OpenAI}}\label{gema-v-openai}

\textbf{Court and date:} Landgericht München I, 42 O 14139/24, 11 Nov
2025

\textbf{What it decided:} Lyrics memorised in model parameters are a
reproduction; TDM exception covers training in general but not
memorisation of specific works; outputs infringe

\textbf{Status on 22 Oct 2026:} Appeal lodged Dec 2025 (OLG München 6 U
3662/25 e), pending

\subsection{\texorpdfstring{\emph{Kneschke v.
LAION}}{Kneschke v. LAION}}\label{kneschke-v-laion}

\textbf{Court and date:} LG Hamburg 27 Sep 2024; OLG Hamburg 10 Dec 2025

\textbf{What it decided:} Downloading images to build a research dataset
is TDM under §44b UrhG (Art. 4 DSM) and also §60d (Art. 3) because LAION
is a non-commercial research organisation; a natural-language
reservation was not shown to be machine-readable

\textbf{Status on 22 Oct 2026:} BGH I ZR 281/25 heard 3 Sep 2026,
decision due 17 Dec 2026, CJEU referral signalled

\subsection{\texorpdfstring{\emph{Like Company v. Google},
C-250/25}{Like Company v. Google, C-250/25}}\label{like-company-v-google-c-25025}

\textbf{Court and date:} CJEU, referred Apr 2025, Grand Chamber hearing
10 Mar 2026

\textbf{What it decided:} Whether LLM training is a reproduction and
whether Art. 4 TDM covers it

\textbf{Status on 22 Oct 2026:} AG opinion not confirmed; judgment not
expected before 2027

\section{9. Pending confirmations}\label{9-pending-confirmations}

UMI puts its characterisation to each rights holder in writing and
records the reply here. A "Pending permission" ruling in section 3
becomes "Allowed" or "Not permitted" when the corresponding line is
updated.

\subsection{\texorpdfstring{Microsoft (ASL STEM Wiki),
ASL\_\href{mailto:Citizen@microsoft.com}{\nolinkurl{Citizen@microsoft.com}}}{Microsoft (ASL STEM Wiki), ASL\_Citizen@microsoft.com}}\label{microsoft-asl-stem-wiki-asl_citizenmicrosoftcom}

\textbf{What UMI has asked:} Written confirmation that UMI cohort
training on v1.0, with CC BY-NC-SA 4.0 public release of weights,
destruction of the dataset and derived files after training, and no
distribution of modifications, is within the Microsoft Research License
Terms and the interpreter consent; or a variation permitting it

\textbf{Status:} Request to be sent; ruling stays Pending permission
until a reply is recorded

\subsection{How2Sign consortium (UPC, BSC, CMU, Gallaudet), via the
CORA.RDR contact and
ELDA}\label{how2sign-consortium-upc-bsc-cmu-gallaudet-via-the-corardr-contact-and-elda}

\textbf{What UMI has asked:} Notice of this policy and of
UMI\textquotesingle s reading that cohort training with CC BY-NC-SA 4.0
release is a non-commercial research use within CC BY-NC 4.0; request
for the CORA.RDR custom dataset terms

\textbf{Status:} Notice to be sent; ruling is Allowed with conditions
and is not contingent on a reply

\subsection{Google Research
(FLEURS-ASL)}\label{google-research-fleurs-asl}

\textbf{What UMI has asked:} Confirmation that release of a mixed model
under CC BY-NC-SA 4.0 with attribution to FLEURS-ASL and FLORES is
acceptable where ShareAlike would otherwise call for CC BY-SA 4.0

\textbf{Status:} Request to be sent; ruling is Allowed and is not
contingent on a reply

\section{Sources}\label{sources}

Dataset terms (opened 22 Oct 2026):
\href{https://www.microsoft.com/en-us/research/project/asl-stem-wiki/}{ASL
STEM Wiki project page},
\href{https://www.microsoft.com/en-us/research/project/asl-stem-wiki/dataset-license/}{Microsoft
Research License Terms},
\href{https://www.microsoft.com/en-us/research/project/asl-stem-wiki/datasheet/}{datasheet},
\href{https://www.microsoft.com/en-us/research/project/asl-stem-wiki/consent-form/}{consent
form},
\href{https://www.microsoft.com/en-us/download/details.aspx?id=106253}{Download
Center entry}; \href{https://how2sign.github.io/}{How2Sign site},
\href{https://arxiv.org/pdf/2008.08143v2}{How2Sign paper},
\href{http://catalog.elra.info/en-us/repository/browse/ELRA-S0416}{ELRA
catalogue ELRA-S0416},
\href{https://www.islrn.org/resources/583-408-694-292-6/}{ISLRN record},
\href{https://b2find.eudat.eu/dataset/99ba11b8-9428-5177-9e25-ccccc296e061}{CORA.RDR
record via B2FIND};
\href{https://www.kaggle.com/datasets/googleai/fleurs-asl}{FLEURS-ASL on
Kaggle} and its Croissant metadata,
\href{https://aclanthology.org/2025.naacl-long.314.pdf}{FLEURS-ASL
paper, NAACL 2025}, \href{https://arxiv.org/pdf/2408.13585}{arXiv
version},
\href{https://huggingface.co/datasets/openlanguagedata/flores_plus}{FLORES+
card}, \href{https://huggingface.co/datasets/google/fleurs}{FLEURS
card}; \href{https://github.com/Umi-BitSign/umi}{UMI subnet repo},
\href{https://github.com/Umi-BitSign/umi-reference-model}{UMI reference
model}.

Licenses and Creative Commons:
\href{https://creativecommons.org/licenses/by-nc/4.0/legalcode}{CC BY-NC
4.0 legal code},
\href{https://creativecommons.org/licenses/by-sa/4.0/legalcode.txt}{CC
BY-SA 4.0 legal code},
\href{https://creativecommons.org/using-cc-licensed-works-for-ai-training-2/}{Using
CC-licensed works for AI training} (28 May 2025),
\href{https://creativecommons.org/faq/\#artificial-intelligence-and-cc-licenses}{CC
FAQ on AI},
\href{https://wiki.creativecommons.org/wiki/NonCommercial_interpretation}{NonCommercial
interpretation},
\href{https://creativecommons.org/2021/12/17/creative-commons-statement-on-cc-licenses-and-the-ext-and-data-mining-exception-under-article-4-eu-cdsm-directive/}{CC
statement on Art. 4 opt-outs} (17 Dec 2021).

Authorities: Bartz v. Anthropic, No. 24-cv-05417 (N.D. Cal. 23 Jun 2025;
final settlement approval 20 Jul 2026); Kadrey v. Meta Platforms (N.D.
Cal. 25 Jun 2025; §1292(b) denied 8 Jul 2026); Thomson Reuters v. Ross
Intelligence, No. 25-2153 (3d Cir. 29 Sep 2026), affirming (D. Del. 11
Feb 2025); US Copyright Office,
\href{https://www.copyright.gov/ai/}{Copyright and Artificial
Intelligence, Part 3} (pre-publication, 9 May 2025); Great Minds v.
FedEx Office, 886 F.3d 91 (2d Cir. 2018); Great Minds v. Office Depot
(9th Cir. 27 Dec 2019);
\href{https://caselaw.findlaw.com/us-federal-circuit/1189790.html}{Jacobsen
v. Katzer}, 535 F.3d 1373 (Fed. Cir. 2008); Zellmer v. Meta Platforms
(9th Cir. 17 Jun 2024);
\href{https://caselaw.nationalarchives.gov.uk/ewhc/ch/2025/2863}{Getty
Images v. Stability AI, {[}2025{]} EWHC 2863 (Ch)} and
\href{https://ipkitten.blogspot.com/2026/01/permission-to-appeal-granted-in-getty.html}{permission
to appeal}; GEMA v. OpenAI, LG München I 42 O 14139/24 (11 Nov 2025),
\href{https://urheber.info/diskurs/openai-legt-berufung-ein}{appeal
lodged},
\href{https://www.morganlewis.com/pubs/2026/06/ai-and-copyright-judicial-landscape-in-germany}{overview};
Kneschke v. LAION,
\href{https://www.bundesgerichtshof.de/SharedDocs/Pressemitteilungen/DE/2026/2026085.html}{BGH
press release on the OLG Hamburg decision} and
\href{https://www.bundesgerichtshof.de/SharedDocs/Termine/DE/Termine/IZR281-25.html}{hearing
listing I ZR 281/25};
\href{https://www.stephensonharwood.com/insights/cjeu-to-rule-on-ai-and-copyright-in-a-landmark-case-against-google/}{Like
Company v. Google, C-250/25}; Directive 2019/790
\href{https://www.legislation.gov.uk/eudr/2019/790/article/2/data.htm}{Art.
2},
\href{https://www.legislation.gov.uk/eudr/2019/790/article/3/data.htm}{Art.
3},
\href{https://www.legislation.gov.uk/eudr/2019/790/article/4/data.htm}{Art.
4};
\href{https://ai-act-service-desk.ec.europa.eu/en/ai-act/article-53}{AI
Act Art. 53};
\href{https://assets.publishing.service.gov.uk/media/69ba692226909a14239612e4/CP2602959_-_Report_on_Copyright_and_Artificial_Intelligence_web.pdf}{UK
Government Report on Copyright and AI, CP 2602959} (18 Mar 2026);
\href{https://www.edpb.europa.eu/our-work-tools/our-documents/opinion-board-art-64/opinion-282024-certain-data-protection-aspects_en}{EDPB
Opinion 28/2024}.

\end{document}
